% ─────────────────────────────────────────────────────────────────────────────
% model_additions.tex
%
% Five derivations that belong in the mathematical model section but were
% either absent or stated incorrectly.  Each is written to be dropped in as a
% subsection of Sec. III; the \label commands match the \ref calls already
% present in the rewritten results.tex.
%
% Insert order:
%   sec:epnp_gn         -> after the EPnP null-space subsection
%   sec:pnp_manifold    -> replaces the Rodrigues-parameterised LM subsection
%   sec:mekf_residual   -> replaces the measurement-model subsection
%   sec:process_noise   -> after the MEKF covariance propagation subsection
%   sec:cone            -> replaces the approach-cone paragraph in the MPC part
% ─────────────────────────────────────────────────────────────────────────────

\subsection{Refinement of the EPnP Control-Point Coefficients}
\label{sec:epnp_gn}

Writing the $N$ world points as affine combinations of control points,
$\mathbf{P}_i = \sum_{j} \alpha_{ij}\,\mathbf{c}_j$ with
$\sum_j \alpha_{ij} = 1$, the projection equations stack into
$\mathbf{M}\mathbf{x} = \mathbf{0}$, where $\mathbf{x}$ collects the control
points in the camera frame. The solution lies in the null space of
$\mathbf{M}$,
\begin{equation}
\mathbf{x} = \sum_{k=1}^{n} \beta_k \mathbf{v}_k ,
\label{eq:epnp_null}
\end{equation}
and the coefficients $\beta_k$ are fixed by requiring that the
inter-control-point distances be preserved:
\begin{equation}
\bigl\| \textstyle\sum_k \beta_k (\mathbf{v}_k^{[i]} - \mathbf{v}_k^{[j]}) \bigr\|^2
= \| \mathbf{c}_i - \mathbf{c}_j \|^2 ,
\qquad i < j .
\label{eq:epnp_dist}
\end{equation}

Equation~\eqref{eq:epnp_dist} is quadratic in $\boldsymbol{\beta}$. Linearising
in the products $b_{kl} = \beta_k \beta_l$ gives a linear system
$\mathbf{L}\mathbf{b} = \boldsymbol{\rho}$ whose least-squares solution yields
an initial $\boldsymbol{\beta}$ by relinearisation,
$\beta_1 = \sqrt{|b_{11}|}$ and $\beta_k = b_{1k}/\beta_1$.

This is an \emph{initialisation}, not the answer, and treating it as the answer
costs roughly a factor of four in accuracy (Sec.~\ref{sec:res_pose}). We refine
$\boldsymbol{\beta}$ by Gauss--Newton directly on \eqref{eq:epnp_dist}. With
$\mathbf{D}_{ij} \in \mathbb{R}^{3 \times n}$ the matrix whose $k$-th column is
$\mathbf{v}_k^{[i]} - \mathbf{v}_k^{[j]}$, and
$\mathbf{d}_{ij}(\boldsymbol{\beta}) = \mathbf{D}_{ij}\boldsymbol{\beta}$, the
residual and its Jacobian are
\begin{align}
f_{ij}(\boldsymbol{\beta}) &= \mathbf{d}_{ij}^{\top}\mathbf{d}_{ij}
                              - \|\mathbf{c}_i - \mathbf{c}_j\|^2 ,
\label{eq:epnp_resid}\\
\frac{\partial f_{ij}}{\partial \beta_m}
  &= 2\, \mathbf{d}_{ij}^{\top} \mathbf{D}_{ij}\,\mathbf{e}_m .
\label{eq:epnp_jac}
\end{align}
Ten iterations of $\boldsymbol{\beta} \leftarrow \boldsymbol{\beta} +
\boldsymbol{\delta}$, with $\boldsymbol{\delta}$ the least-squares solution of
$\mathbf{J}\boldsymbol{\delta} = -\mathbf{f}$, suffice in every case tested.

\paragraph{The coplanar case.}
The control points are the centroid plus the principal axes of the point
cloud, scaled by the corresponding singular values. When the cloud is
coplanar the third singular value vanishes, so a fourth control point
coincides with the centroid, the barycentric system becomes singular, and
three columns of $\mathbf{M}$ vanish identically --- giving $\mathbf{M}$ a
four-dimensional null space for algebraic rather than geometric reasons. The
planar configuration must therefore be detected, from
$s_3 / s_1 < \epsilon$, and parameterised with \emph{three} control points, so
that $\mathbf{x} \in \mathbb{R}^{9}$ and $\mathbf{M}$ is $2N \times 9$. This
is not an edge case: approximately 3\% of random six-point subsets of the
target's keypoint set are exactly coplanar.

\subsection{Pose Refinement on the $SO(3)$ Manifold}
\label{sec:pnp_manifold}

The refinement stage minimises the reprojection cost
\begin{equation}
E(\mathbf{R}, \mathbf{t}) = \sum_{i} w_i \,
\bigl\| \tilde{\mathbf{p}}_i - \pi(\mathbf{R}\mathbf{P}_i + \mathbf{t}) \bigr\|^2 .
\label{eq:lm_cost}
\end{equation}
Parameterising $\mathbf{R} = \exp([\mathbf{r}]_\times)$ by a \emph{global}
rotation vector is convenient but singular at $\|\mathbf{r}\| = \pi$, where the
derivative of the exponential map is unbounded. This is not a remote
possibility here: a camera whose boresight is aligned with a world axis
produces $\operatorname{tr}(\mathbf{R}) = -1$ exactly, i.e.\
$\theta = \pi$, and such poses are the natural fixtures for an on-axis
approach.

We therefore carry $\mathbf{R}$ as a matrix and optimise over a local
increment. Using a left-multiplicative update,
\begin{equation}
\mathbf{R} \leftarrow \exp([\boldsymbol{\delta\omega}]_\times)\,\mathbf{R},
\qquad
\mathbf{t} \leftarrow \mathbf{t} + \boldsymbol{\delta t} ,
\label{eq:lm_update}
\end{equation}
and noting that $\exp([\boldsymbol{\delta}]_\times)\mathbf{v} \approx
\mathbf{v} + \boldsymbol{\delta} \times \mathbf{v} = \mathbf{v} -
[\mathbf{v}]_\times \boldsymbol{\delta}$, the Jacobian blocks at
$\boldsymbol{\delta} = \mathbf{0}$ are
\begin{equation}
\frac{\partial \mathbf{P}^c_i}{\partial \boldsymbol{\delta\omega}}
  = -[\mathbf{R}\mathbf{P}_i]_\times ,
\qquad
\frac{\partial \mathbf{P}^c_i}{\partial \boldsymbol{\delta t}} = \mathbf{I}_3 ,
\label{eq:lm_jac}
\end{equation}
so that with the perspective Jacobian
\begin{equation}
\frac{\partial \pi}{\partial \mathbf{P}^c}
= \begin{bmatrix}
f_x / Z & 0 & -f_x X / Z^2 \\
0 & f_y / Z & -f_y Y / Z^2
\end{bmatrix}
\label{eq:proj_jac}
\end{equation}
the stacked Jacobian is
$\mathbf{J}_i = (\partial \pi / \partial \mathbf{P}^c)\,
[-[\mathbf{R}\mathbf{P}_i]_\times \;\; \mathbf{I}_3]$.
Because $\boldsymbol{\delta\omega}$ is always near zero, the local
parameterisation is exact and well conditioned wherever the optimiser is, and
the singularity cannot be reached.

Two further conditions keep the iteration honest. The residual and the
Jacobian must be evaluated at the same point --- clamping $Z$ in one and not
the other lets Levenberg--Marquardt accept steps that lower a cost it is not
actually computing, and the translation then diverges without bound while the
reported cost stays finite. And chirality must be enforced: any pose placing a
weighted point at or behind the camera plane is assigned infinite cost and
rejected, so failure is visible rather than silent.

\subsection{Multiplicative Attitude Residual}
\label{sec:mekf_residual}

The measurement supplies a relative position and an attitude. Position enters
as a difference in the usual way. Attitude cannot, and the reason is worth
setting out because the naive form is superficially reasonable.

Let $\mathbf{q}$ denote the attitude quaternion and $\hat{\mathbf{q}}$ the
filter's estimate. Encoding both as global rotation vectors and subtracting,
$\boldsymbol{\nu}_{att} = \log(\mathbf{q}_{meas}) - \log(\hat{\mathbf{q}})$,
fails at $\theta = \pi$: two attitudes separated by an arbitrarily small
rotation, one either side of $\pi$, map to nearly antipodal rotation vectors,
so a vanishing attitude error produces a residual of magnitude $2\pi$. The
linearised residual is then meaningless, and the Jacobian condition number
grows without bound as $\theta \to \pi$.

The residual must instead be formed on the group and expressed in the tangent
space at the estimate:
\begin{equation}
\boldsymbol{\delta\alpha}
  = \log\!\bigl( \mathbf{q}_{meas} \otimes \hat{\mathbf{q}}^{-1} \bigr) ,
\label{eq:att_residual}
\end{equation}
where $\log(\cdot)$ returns the rotation vector of its argument and the sign of
the relative quaternion is chosen so that its scalar part is non-negative,
selecting the short rotation. This is exact for every attitude.

Equation~\eqref{eq:att_residual} has a second benefit. Because the error state
is defined by the same left-multiplicative perturbation,
$\mathbf{q} = \exp(\tfrac{1}{2}\boldsymbol{\delta\alpha}) \otimes
\hat{\mathbf{q}}$, the residual \emph{is} the attitude error state to first
order, and the measurement Jacobian becomes exact and constant:
\begin{equation}
\mathbf{H} =
\begin{bmatrix}
\mathbf{I}_3 & \mathbf{0} & \mathbf{0} & \mathbf{0} \\
\mathbf{0} & \mathbf{0} & \mathbf{I}_3 & \mathbf{0}
\end{bmatrix} .
\label{eq:H_exact}
\end{equation}
No finite differencing is required, and $\mathbf{R}_{att} = \sigma_{att}^2
\mathbf{I}_3$ becomes a genuine tangent-space covariance rather than an
approximation. The measurement noise must be injected on the same side for
this to hold: a physical attitude perturbation is
$\mathbf{q}_{meas} = \exp(\tfrac{1}{2}\boldsymbol{\eta}) \otimes \mathbf{q}$
with $\boldsymbol{\eta} \sim \mathcal{N}(\mathbf{0},
\sigma_{att}^2\mathbf{I}_3)$.

\subsection{Acceleration-Driven Process Noise}
\label{sec:process_noise}

The dominant unmodelled effects on the relative translation --- differential
$J_2$, differential drag, thruster execution error --- act as an
\emph{acceleration}. A process noise model must therefore drive position and
velocity jointly. Treating the disturbance as white acceleration of density
$\sigma_a$ over a step $\Delta t$ gives the standard kinematic block
\begin{equation}
\begin{bmatrix}
\mathbf{Q}_{rr} & \mathbf{Q}_{rv} \\
\mathbf{Q}_{vr} & \mathbf{Q}_{vv}
\end{bmatrix}
= \sigma_a^2
\begin{bmatrix}
\tfrac{1}{3}\Delta t^3 \mathbf{I}_3 & \tfrac{1}{2}\Delta t^2 \mathbf{I}_3 \\
\tfrac{1}{2}\Delta t^2 \mathbf{I}_3 & \Delta t\, \mathbf{I}_3
\end{bmatrix},
\label{eq:Q_kin}
\end{equation}
with the same structure on the attitude and rate block for an angular
acceleration density.

The off-diagonal terms are not optional. Position and velocity errors driven by
a common acceleration are correlated, and a diagonal $\mathbf{Q}$ asserts that
they are not, which weakens the information that a position measurement carries
about velocity. A diagonal alternative that models position as an independent
random walk is worse still: a position noise of $0.05$~m per step contributes
$\mathbf{Q}_{rr} = 2.5 \times 10^{-3}$~m$^2$, roughly three hundred times the
$8.3 \times 10^{-6}$~m$^2$ that a physically plausible
$5 \times 10^{-4}$~m/s$^2$ disturbance produces over a one-second step. The
filter is then needlessly conservative, and --- more importantly for
validation --- its $\mathbf{Q}$ describes a disturbance that is not present in
the truth, so the consistency statistics of Eq.~\eqref{eq:nis} become
uninformative.

\subsection{Approach Corridor as a Second-Order Cone}
\label{sec:cone}

The terminal approach is confined to a cone of half-angle $\theta_c$ about the
$+x$ (radial) approach axis. For a predicted position
$\mathbf{r}_k = [x_k,\ y_k,\ z_k]^{\top}$ the constraint is
\begin{equation}
\sqrt{y_k^2 + z_k^2} \;\le\; x_k \tan\theta_c ,
\qquad k = 0,\dots,N-1 .
\label{eq:cone}
\end{equation}
Three properties of \eqref{eq:cone} matter and are each easy to lose.

It is a \emph{cone}, not a wedge: both transverse components appear.
Constraining $|y_k|$ alone leaves the cross-track axis entirely free. It is
\emph{signed}: since the right-hand side must be non-negative,
\eqref{eq:cone} implies $x_k \ge 0$, whereas replacing $x_k$ by $|x_k|$
mirrors the corridor into $x < 0$ and is satisfied by a chaser sitting behind
the target. And the constrained quantity is the predicted state
$\mathbf{X}(\mathbf{U}) = \mathbf{S}_x \mathbf{e} + \mathbf{S}_u \mathbf{U}$;
evaluating the radius on the free response $\mathbf{S}_x\mathbf{e}$ makes the
right-hand side independent of the decision variable, and the constraint
ceases to constrain anything.

Written for a nonlinear-programming solver, \eqref{eq:cone} becomes
$g_k(\mathbf{U}) \ge 0$ with
\begin{equation}
g_k(\mathbf{U}) = \tan\theta_c \; x_k(\mathbf{U})
   - \sqrt{y_k(\mathbf{U})^2 + z_k(\mathbf{U})^2 + \varepsilon^2},
\label{eq:cone_g}
\end{equation}
where $\varepsilon$ (here $10^{-6}$~m) keeps $g_k$ differentiable on the axis
at the cost of a marginally conservative constraint. The gradient is
\begin{equation}
\nabla_{\mathbf{U}} g_k = \tan\theta_c \, \mathbf{S}_u^{x,k}
 - \frac{y_k \mathbf{S}_u^{y,k} + z_k \mathbf{S}_u^{z,k}}
        {\sqrt{y_k^2 + z_k^2 + \varepsilon^2}} ,
\label{eq:cone_grad}
\end{equation}
with $\mathbf{S}_u^{x,k}$ the row of $\mathbf{S}_u$ selecting $x_k$.

\paragraph{Recursive feasibility.}
Imposed as a hard constraint, \eqref{eq:cone} destroys recursive feasibility:
once the chaser lies outside the cone there may be no admissible input
returning the first predicted state to it, and the quadratic program is simply
infeasible. Since an infeasible solve is what produces a silent
zero-thrust command, we soften the corridor with slack
$s_k \ge 0$ under a linear penalty,
\begin{equation}
\min_{\mathbf{U},\,\mathbf{s}} \;
\tfrac{1}{2}\mathbf{U}^{\top}\mathbf{H}\mathbf{U} + \mathbf{f}^{\top}\mathbf{U}
+ \rho \sum_{k} s_k
\quad \text{s.t.} \quad g_k(\mathbf{U}) + s_k \ge 0 ,
\label{eq:cone_soft}
\end{equation}
which is always feasible and drives $s_k \to 0$ whenever the hard constraint
can be met. The corridor solve is judged on the quantity it exists to reduce,
$\sum_k \max(0, -g_k)$, rather than on the solver's own convergence flag, and
is attempted only when the box-constrained optimum already violates
\eqref{eq:cone} --- an active-set shortcut that leaves the solution unchanged
when the corridor is inactive and reduces the per-step cost by two orders of
magnitude.
